% Source: https://artofproblemsolving.com/community/c3h1772123_mock_target_8__factorial_number_bases
% Stable identity: aops-topic-1772123-q01
Every positive integer has exactly one base factorial representation (the factorial number system). Since $1\cdot1!+2\cdot2!+\cdots+k\cdot k!=(k+1)!-1$, the numbers with exactly $k$ digits are those from $k!$ to $(k+1)!-1$. So the numbers $1$ to $719$ have at most $5$ digits, and $720=6!$ is $(1000000)$, which contains a $0$.

Work with the complement: the numbers whose digits are all nonzero. For a $k$-digit such number, each digit $d_i$ ranges over $1,2,\dots,i$ independently, giving $k!$ numbers, and over these numbers each $d_i$ takes each value equally often, with average $\dfrac{i+1}{2}$. So their sum is
\[
k!\sum_{i=1}^{k}\dfrac{i+1}{2}\,i!=\dfrac{k!}{2}\sum_{i=1}^{k}(i+1)! .
\]
For $k=1,\dots,5$ this gives $1$, $8$, $3\cdot32=96$, $12\cdot152=1824$ and $60\cdot872=52320$, a total of $54249$. Hence
\[
N=\dfrac{720\cdot721}{2}-54249=259560-54249=205311 .
\]
Since $205311=285\cdot720+111$, the remainder is $\boxed{111}$.
